\notechapter{N-20221221}{Basic Topology II: Completion, Homeomorphism Warnings, and Quotient Intuition}

\section{Completion}

The note continues with the theorem that every metric space can be completed.  The standard example is
\[
  \mathbb Q\quad\leadsto\quad\mathbb R.
\]
A rational Cauchy sequence should be thought of as trying to name a real number.  If the real number is irrational, the sequence has no limit in \(\mathbb Q\), but it has a limit after completion.

One construction takes equivalence classes of Cauchy sequences.  Two Cauchy sequences \((x_n)\), \((y_n)\) represent the same completed point if
\[
  d(x_n,y_n)\to0.
\]
This is an early quotient construction: the new point is not one sequence but a class of mutually converging sequences.

\section{Let the buyer beware}

The note warns that boundedness and completeness are not topological invariants.  The interval \((0,1)\) and the line \(\mathbb R\) are homeomorphic, for instance by
\[
  f(x)=\tan(\pi(x-1/2)).
\]
But \((0,1)\) is bounded and not complete under the Euclidean metric, while \(\mathbb R\) is complete and unbounded.

This is the first serious conceptual warning: homeomorphism preserves open-set structure, not distances, diameters, or Cauchy behavior.

\section{Homeomorphisms and embeddings}

A homeomorphism \(f:M\to N\) is a bijection such that both \(f\) and \(f^{-1}\) are continuous.  The inverse condition is essential.  A continuous bijection can fail to be a homeomorphism if the inverse destroys neighborhoods.

An embedding realizes a space as a subspace of another space in a topology-preserving way.  It is more than an injection: it must preserve the local structure inherited from the target.

\section{Quotient intuition}

A quotient space identifies points.  For example,
\[
  [-1,1]/(-1\sim1)\cong S^1.
\]
A neighborhood of the glued point corresponds upstairs to two half-neighborhoods at the two endpoints.  This explains why quotient topology is defined by requiring a set downstairs to be open exactly when its preimage upstairs is open.

\section{A more structural view}

Completion adds missing limits.  Homeomorphism forgets metric size.  Embedding preserves a space as a subspace.  Quotient topology glues points together.  These are four ways of changing a space while controlling which structure survives.

\section{Completion and its universal property}

The completion \(X\hookrightarrow \widehat X\) has a useful universal property.  If \(Y\) is complete and \(f:X\to Y\) is uniformly continuous, then \(f\) extends uniquely to a continuous map
\[
  \widehat f:\widehat X\to Y.
\]
On a completed point \([(x_n)]\), define
\[
  \widehat f([(x_n)])=\lim_{n\to\infty}f(x_n).
\]
Uniform continuity ensures that \((f(x_n))\) is Cauchy in \(Y\), and completeness of \(Y\) supplies the limit.

This explains why completion is canonical.  It is not merely "adding some missing points"; it is the best possible way to make Cauchy limits available without changing the original metric information.

\section{Homeomorphism versus uniform equivalence}

The example \((0,1)\cong\mathbb R\) shows that topology does not preserve boundedness or completeness.  A sharper statement is that the map is a homeomorphism but not a uniform equivalence for the usual metrics.

A uniformly continuous map controls all small distances by one global \(\delta\).  The tangent map
\[
  x\mapsto \tan(\pi(x-1/2))
\]
blows up near the endpoints.  Points very close to \(1\) can be very close in \((0,1)\) but sent far apart in \(\mathbb R\).  Thus topology sees local neighborhoods, while uniform structure sees global control of closeness.

\section{Continuous bijections that are not homeomorphisms}

A continuous bijection need not be a homeomorphism.  A standard example is
\[
  f:[0,2\pi)\to S^1,\qquad f(t)=(\cos t,\sin t).
\]
It is continuous and bijective, but its inverse is not continuous at the point \((1,0)\).  Points on the circle approaching \((1,0)\) from below have parameters approaching \(2\pi\), while the point itself has parameter \(0\).

This example explains why the inverse condition in the definition of homeomorphism is not pedantry.  A bijection can preserve points but fail to preserve neighborhoods.

\section{Quotient topology as final topology}

For a quotient map \(q:X\to Y\), the quotient topology on \(Y\) is the strongest topology making \(q\) continuous.  It is characterized by
\[
  U\subseteq Y\text{ open}
  \quad\Longleftrightarrow\quad
  q^{-1}(U)\text{ open in }X.
\]
The universal property is:
\[
  f:Y\to Z\text{ is continuous}
  \quad\Longleftrightarrow\quad
  f\circ q:X\to Z\text{ is continuous}.
\]
This is the correct abstract version of endpoint gluing.  To define a continuous function out of a glued space, define it before gluing and check that it is constant on equivalence classes.

\section{Endpoint gluing and local neighborhoods}

In \([0,1]/(0\sim1)\), a neighborhood of the glued point is not a one-sided interval.  Its preimage is a union of a small interval near \(0\) and a small interval near \(1\):
\[
  [0,\varepsilon)\cup(1-\varepsilon,1].
\]
This two-sided local structure is why the quotient is a circle.  The topology is doing exactly what the picture suggests: it forces the two endpoints to share neighborhoods after identification.

\section{Completion, gluing, and quotient control}

Completion, homeomorphism, embedding, and quotient are not interchangeable changes of space.  Completion preserves metric Cauchy information while adding limits.  Homeomorphism preserves open-set structure and may forget metric size.  Embedding realizes a topology inside another space.  Quotient topology glues points and controls maps out of the glued space.
